\section{Results}
\label{sec:results}

\subsection{Main Out-of-Sample Results}
\label{subsec:results_main}

Table~\ref{tab:main_results} reports out-of-sample performance for every
strategy on the common evaluation window described in
Section~\ref{subsec:walkforward}: the 720 trading days between 2022-01-03 and
2024-11-13 on which the agent and all three benchmarks produce a return.
The four agent configurations are trained under an identical protocol
(5{,}000 gradient steps, seed 42) and differ only in the component named in
Section~\ref{subsec:ablation}.
The metrics reported are annualized return, Sharpe ratio, Sortino ratio,
CVaR$_{0.95}$ (defined as the average loss in the worst 5\% of daily return
observations, reported as a positive number for clarity), maximum drawdown (MaxDD),
and Calmar ratio (annualized return divided by absolute maximum drawdown).
Every figure comes from a single evaluation window and a single seed and carries
no interval estimate; Section~\ref{subsec:disc_limitations} sets out what that
precludes.

Among the three classical benchmarks, Equal-Weight delivers the highest Sharpe
ratio (0.6076) and annualized return (10.34\%), confirming the $1/N$ puzzle
documented by \citet{DeMiguel2007} in the context of B3 equities.
The Minimum-Variance Portfolio achieves the lowest CVaR$_{0.95}$ (0.0198) and the
smallest maximum drawdown ($-18.22\%$) but at the cost of markedly lower returns
(5.16\% annualized), yielding a Sharpe ratio of 0.3400.
Risk Parity occupies an intermediate position, with a Sharpe ratio of 0.5561 and
an annualized return of 8.94\%.

The full CQL agent (configuration A4) achieves a Sharpe ratio of 0.4058 and an
annualized return of 7.88\%, with a CVaR$_{0.95}$ of 0.0260 and a maximum
drawdown of $-27.22\%$.
This is $-0.2018$ in Sharpe relative to Equal-Weight and $+0.0657$ relative to
MVP, the strategy that generated the behavioral data.
Equal-Weight dominates A4 on every metric in Table~\ref{tab:main_results}:
higher return, higher Sharpe, higher Sortino, lower tail loss, shallower drawdown
and higher Calmar.
Risk Parity likewise dominates it.
Against MVP the picture is mixed rather than favourable: A4 earns more and has a
higher Sharpe, but its tail loss is a third larger (0.0260 against 0.0198) and its
drawdown half again as deep.

Two features of the table deserve to be stated plainly, because the rest of this
section is an attempt to account for them.
First, the proposed configuration is not the best agent in its own ablation: A2,
which uses neither the CVaR reward nor the macro state, attains a Sharpe of 0.4464
against A4's 0.4058.
Second, CVaR$_{0.95}$ increases monotonically as the CVaR reward is introduced and
then extended: 0.0222 for A1, 0.0221 for A2, 0.0241 for A3 and 0.0260 for A4.
The two configurations that optimize a tail-risk objective have the worst realised
tail loss of the seven strategies compared.

\inputtable{main_results}

\begin{figure}[htbp]
  \centering
  \includegraphics[width=\linewidth]{fig_wealth_curves}
  \caption{Cumulative wealth index for all strategies, out-of-sample period 2022--2024
           (base = 1 at start of evaluation).
           A4 (CQL+CVaR+Macro) is plotted with a thicker line.}
  \label{fig:wealth_curves}
\end{figure}

\begin{figure}[htbp]
  \centering
  \includegraphics[width=\linewidth]{fig_drawdown}
  \caption{Drawdown trajectories for A4, Equal-Weight, MVP, and Risk Parity,
           out-of-sample period 2022--2024.}
  \label{fig:drawdown}
\end{figure}

\subsection{Sub-Period Analysis}
\label{subsec:results_subperiod}

We divide the out-of-sample period into three sub-periods to assess performance
stability across distinct macroeconomic regimes:

\textbf{2022 tightening phase.}
This period corresponds to the most aggressive phase of the SELIC hiking cycle,
during which the rate rose from 9.25\% to 13.75\%.
The agent does not provide downside protection here.
A4 returns $-1.94\%$ annualized with a Sharpe of $-0.086$ and a
CVaR$_{0.95}$ of 0.0297, against $+11.89\%$, 0.573 and 0.0265 for Equal-Weight.
This is the sub-period in which the motivating hypothesis predicted the clearest
advantage, and it is the sub-period in which the agent performs worst in both
return and tail loss.

\textbf{2023 peak rate period.}
With the SELIC at its ceiling of 13.75\%, A4 generates a Sharpe of 0.8374 against
1.3501 for Equal-Weight, and a larger tail loss (0.0243 against 0.0200).
The agent is profitable in absolute terms and still behind the naive benchmark.

\textbf{2024 easing phase.}
This is the one sub-period in which A4 wins, and it wins decisively: 9.50\%
annualized with a Sharpe of 0.633, against $-5.05\%$ and $-0.424$ for
Equal-Weight.
Taken together with the two preceding sub-periods, the pattern is not one of
regime intelligence but of persistent defensive positioning: the agent gives up
return when equities rise and preserves capital when they fall.
A lower-beta portfolio would produce the same three-period signature without any
learning at all, and this study does not include the beta-matched comparison that
would separate the two explanations.

\subsection{Regime-Conditional Performance}
\label{subsec:results_regime}

A central hypothesis of this paper is that macro regime information improves
performance specifically during adverse monetary-policy environments.
To test this, we compute performance metrics conditional on the SELIC regime
label at each trading day.

Table~\ref{tab:regime_conditional} reports Sharpe ratios and CVaR by regime.
The hypothesis is not supported, and the ordering it predicts is reversed.

In the Rising-rate regime, A4 achieves a Sharpe of 0.2811 against 1.2050 for
Equal-Weight: the tightening phase is where the agent trails the naive benchmark
by the widest margin, not the narrowest.
In the Stable regime the gap narrows but persists, 0.4929 against 0.7318.
In the Falling regime the ranking inverts: A4 achieves 0.4609 while Equal-Weight
returns $-0.4524$ and MVP $-0.2184$, making the agent the only strategy of the
three with a positive Sharpe while rates come down.

Note that the sign of the market's own regime-conditional performance differs
between the full sample and the evaluation window.
Over 2015--2024 the equal-weight portfolio earns an annualized Sharpe of $-0.23$
in rising-rate phases (Table~\ref{tab:regime_stats}), which is the empirical
regularity that motivated conditioning on the cycle in the first place.
Within the 2022--2024 evaluation window that regularity does not hold: equities
rose through the tightening phase and fell through the easing phase.
An agent whose training data encodes the historical association therefore
inherits a mapping that the test period does not honour, and the regime-conditional
results are as much a statement about that instability as about the agent.

\inputtable{regime_conditional}

\inputtable{subperiod_results}

\begin{figure}[htbp]
  \centering
  \includegraphics[width=\linewidth]{fig_regime_sharpe}
  \caption{Regime-conditional Sharpe ratios for A4 and Equal-Weight
           across Rising, Stable, and Falling SELIC regimes.}
  \label{fig:regime_sharpe}
\end{figure}

\begin{figure}[htbp]
  \centering
  \includegraphics[width=0.8\linewidth]{fig_cvar_subperiods}
  \caption{CVaR$_{0.95}$ by sub-period for A4, Equal-Weight, and MVP.
           Lower values indicate better tail-risk control.}
  \label{fig:cvar_subperiods}
\end{figure}

\subsection{Ablation Study}
\label{subsec:results_ablation}

Table~\ref{tab:ablation_results} reports the four configurations A1--A4.
Performance does not improve monotonically along the ladder, and the two
components the paper set out to defend are the two that cost performance.

\textbf{A1 vs.\ A2: the conservatism penalty.}
Raising the conservative weight from $0.001$ to the value selected by grid search
lifts Sharpe from 0.1557 to 0.4464, leaves tail loss essentially unchanged
(0.0222 to 0.0221) and increases maximum drawdown modestly ($-23.14\%$ to
$-25.49\%$ before rounding).
This is the largest single gain in the study and the only component that survives
the ablation.
It is also the result most consistent with the offline RL literature: penalising
value estimates for actions absent from the behavioral data is what keeps an
offline agent from extrapolating into regions the dataset cannot support.
Note that this comparison says nothing about online versus offline training.
Both configurations are offline; only the penalty differs.

\textbf{A2 vs.\ A3: the CVaR reward.}
Replacing the realised-return reward with the CVaR objective lowers Sharpe from
0.4464 to 0.2111 and raises CVaR$_{0.95}$ from 0.0221 to 0.0241.
The reward intended to control tail risk makes tail risk worse.
Section~\ref{subsec:disc_cvar} argues that this follows from how the signal is
constructed rather than from any property of the market: the reward scores
candidate weights against the same twenty-day window the state already encodes,
so it carries no information the agent does not already hold and provides no
feedback about the consequences of the action.

\textbf{A3 vs.\ A4: the macro state.}
Adding the SELIC target, the lending spread and the regime encoding lifts Sharpe
from 0.2111 to 0.4058 while raising CVaR$_{0.95}$ further, from 0.0241 to 0.0260,
and deepening maximum drawdown to $-27.22\%$.
The macro block recovers roughly four fifths of the Sharpe lost to the reward
change without recovering any of the tail deterioration.
Two readings are available and the data does not separate them.
The macro variables may carry genuine information about the allocation decision;
or, because the reward supplies no feedback, the additional state dimensions may
simply be changing which trailing-window portfolio the agent converges to.
Establishing the first reading requires re-running A3 and A4 under a reward that
is not measurable from the state, which is the experiment
Section~\ref{subsec:disc_cvar} identifies as necessary.

\textbf{The ladder as a whole.}
The best agent in the study is A2, not the proposed A4, and A2 uses neither of the
two components this paper contributes.
Every agent configuration is below Equal-Weight and below Risk Parity on Sharpe.
The honest summary is that of the three ingredients tested, one helps, one hurts,
and the third cannot be interpreted until the second is fixed.

\inputtable{ablation_results}

\begin{figure}[htbp]
  \centering
  \includegraphics[width=\linewidth]{fig_ablation_heatmap}
  \caption{Ablation heatmap: normalized performance metrics for configurations A1--A4.
           Green (1.0) indicates the best value in each metric column;
           red (0.0) indicates the worst.
           Raw metric values are shown in each cell.}
  \label{fig:ablation_heatmap}
\end{figure}
